\section*{Chapitre \ref{ch:rs:event-driven-backtests} --- Backtests événementiels}

\begin{solution}{exo:rs:event-driven-backtests:1}
\omterm{def:rs:event-driven-backtests:fill}{Exécution au contact}~: le plus bas a atteint 99.00, donc oui, en totalité à 99.00. Pénétration d'un tick~: le plus bas aurait dû atteindre 98.99~; ce n'est pas le cas, donc non.
\end{solution}

\begin{solution}{exo:rs:event-driven-backtests:2}
2.5~\% de chaque barre~: 1\,000, 250 et 1\,500 actions, soit 2\,750 au total, laissant 2\,250 en attente après la troisième barre.
\end{solution}

\begin{solution}{exo:rs:event-driven-backtests:3}
en attente (envoyé) $\to$ actif (accusé de réception) $\to$ partiel (100 exécutés, 200 restants) $\to$ annulé (les 200 retirés lorsque l'annulation prend effet)~;
les 100 actions exécutées restent exécutées.
\end{solution}

\begin{solution}{exo:rs:event-driven-backtests:4}
Une transaction au-delà du prix de l'ordre signifie que toute la file à ce prix a été consommée, ce qui se produit lorsque le flux de l'autre côté est important et
persistant, souvent informé (chapitre~9)~; le prix continue alors de bouger. Un simple contact peut survenir alors que la file est à peine entamée et que le prix repart en arrière. La
pénétration d'exécution sélectionne donc les exécutions qu'obtiendrait un ordre réel placé en queue de file, et ce sont les exécutions défavorables~: sur les journées synthétiques,
un mark-out de $-2.0$ ticks contre $-0.16$.
\end{solution}

\begin{solution}{exo:rs:event-driven-backtests:5}
Chaque ordre est envoyé à la clôture d'une barre, devient actif 5 secondes après le début de la barre suivante et est annulé 5 secondes après le début de celle d'après~; il n'est jamais actif
pendant une barre entière. Une stratégie qui laisse ses ordres en place tant que le prix ne change pas les garde actifs pendant de nombreuses barres entières, et la politique conservatrice peut les exécuter.
\end{solution}

\begin{solution}{exo:rs:event-driven-backtests:6}
À l'instant d'effet de la division~: position portée à 200 actions, ordre porté à 200 actions à 60~\$. Sans cela, l'ordre demande encore 120~\$ pour 100 actions dans un marché
désormais autour de 60~\$~: il ne s'exécute jamais (le livre reste sans couverture face à la sortie prévue), ou, pour un achat à cours limité laissé au prix ancien, il s'exécute aussitôt au double
du marché.
\end{solution}

\begin{solution}{exo:rs:event-driven-backtests:7}
\texttt{rs\_evbt.at\_close}~: avec l'\omterm{def:rs:event-driven-backtests:fill}{exécution au contact}, 95~\% des ordres s'exécutent, avec un mark-out de $-0.17$ ticks et une perte de 188~\$ par jour~; avec l'\omterm{def:rs:event-driven-backtests:fill}{exécution par pénétration}
70~\% s'exécutent, avec $-1.16$ ticks et une perte de 506~\$ par jour. Les ordres au prix qui vient d'être traité sont presque toujours touchés, ce qui ne dit rien de ce qu'il en serait pour un
ordre réel placé en queue de file serait exécuté.
\end{solution}

\begin{solution}{exo:rs:event-driven-backtests:8}
Le \omterm{def:rs:event-driven-backtests:fill}{modèle d'exécution} suppose que l'ordre était en tête de file à son prix et l'exécute à chaque contact, y compris à tous les contacts où le prix repart en
arrière, c'est-à-dire les exécutions qu'un ordre réel n'obtiendrait pas~; il ignore la sélection des exécutions défavorables. Il faut relancer avec des \omterm{def:rs:event-driven-backtests:fill}{exécutions par pénétration} et un plafond de volume, vérifier les mark-outs,
et confirmer au niveau~3 ou avec des exécutions réelles avant de lui faire confiance.
\end{solution}

\begin{solution}{pb:rs:event-driven-backtests:1}
\begin{enumerate}
\item Exécutions contre la barre d'abord (les ordres étaient actifs pendant celle-ci), puis marquage à sa clôture, puis réaction de la stratégie~; ainsi la
stratégie voit les exécutions et le portefeuille qu'elle détenait réellement, et ses nouveaux ordres n'agissent qu'à partir des barres suivantes.
\item Lorsque son événement d'accusé de réception arrive, après la latence de saisie des ordres~; et, selon la politique conservatrice, seulement dans les barres qui commencent après cet instant.
\item Un ordre total sur les événements (temps, nature, numéro de séquence), un aléa initialisé et transmis explicitement, aucune dépendance à l'horloge murale.
\item Le \texttt{BacktestResult} du chapitre~16, les ordres avec leurs cycles de vie, et les exécutions.
\item Contact~: 523 lots, 70~\%, $+\$17$, $-0.16$ ticks. Pénétration~: 290 lots, 41~\%, $-\$584$, $-2.00$ ticks. Plafonné à 2.5~\%~: 388 lots, 55~\%, $+\$20$,
$-0.18$ ticks.
\item Contact~: en tête de file. Pénétration~: en queue de file. Plafonné~: une part fixe du volume, quelle que soit la position de l'ordre dans la file.
\item Ses exécutions ne sont que celles où le prix a traversé, lesquelles sont suivies de mouvements contre la position~; l'\omterm{def:rs:event-driven-backtests:fill}{exécution au contact} ajoute les exécutions bénignes
où le prix repart en arrière.
\item Détenir un lot en position longue après une barre baissière et un lot en position courte après une barre haussière, échangés à la clôture avec un demi-tick de coût~: $-\$215$ par jour.
\item Aucune exécution~: aucun ordre n'est actif pendant une barre entière.
\item Contact~: 585 lots, $-\$41$~; pénétration~: 372 lots, $-\$776$~; plafonné~: 440 lots, $-\$190$, avec un mark-out qui tombe à $-0.87$ ticks.
\item Le fait que la réponse dépende de l'instant, à l'intérieur de la barre, où le prix a atteint les ordres, que les barres n'enregistrent pas~: la stratégie doit être jugée au niveau~3.
\item Maintenir les ordres actifs tant que le prix ne change pas, n'annuler qu'en cas de mouvement, et coter plus loin du prix.
\item \textbf{Résultat nommé.} Sans latence~: contact, 70~\% des ordres exécutés et $+\$17$ par jour~; pénétration, 41~\% et $-\$584$~; plafonné, 55~\% et $+\$20$~; avec une
latence de 5 secondes~: rien selon la politique conservatrice, et $-\$41$, $-\$776$ et $-\$190$ selon la politique optimiste. Le \omterm{def:rs:event-driven-backtests:fill}{modèle d'exécution} déplace la réponse de 800~\$
par jour~; l'avantage propre de la stratégie reste dans son bruit.
\item Le résultat par pénétration, retenu comme borne pessimiste, avec l'affirmation que le résultat par contact suppose une place en tête de chaque file et que les résultats au
niveau~2 encadrent la réponse sans la localiser.
\item Nos propres ordres dans la file à leurs prix, chaque transaction à ces prix, et la latence côté ordres comme côté données~: des exécutions qui ont lieu ou non
pour les raisons qu'elles auraient dans la réalité.
\item Au plus 2.5~\% de chaque barre s'exécute~; une stratégie qui en exige 5~\% laisse la moitié de ses ordres actifs et négocie plus tard, à d'autres prix.
\item Les divisions, les dividendes (pour la trésorerie et l'écart de cours), les demi-journées, les enchères et les suspensions de cotation.
\item Parce qu'un ordre réel rejoint l'arrière d'une file qui peut contenir des milliers d'actions à son prix~; pour s'exécuter 70~\% du temps, le prix devrait
traverser cette file pendant la plupart des minutes.
\item L'envoi, l'accusé de réception, l'estimation de la file, les exécutions et les annulations de chaque ordre, horodatés, ainsi que les transactions du marché à son prix~; le rapport entre les
exécutions réelles et les \omterm{def:rs:event-driven-backtests:fill}{exécutions au contact} étalonne le modèle.
\item L'hypothèse du \omterm{def:rs:vectorised-backtests:backtest}{backtest} sur la place de ses ordres dans la file, rendue explicite et remplaçable.
\end{enumerate}
\end{solution}

\begin{solution}{iq:rs:event-driven-backtests:1}
Une file de priorité d'événements horodatés (données de marché, accusés de réception, exécutions, annulations, événements de calendrier et d'entreprise), un répartiteur qui les
extrait dans l'ordre et fait avancer l'\omterm{def:rs:event-driven-backtests:evbt}{horloge simulée}, des gestionnaires pour la stratégie, le gestionnaire d'ordres (cycle de vie), le \omterm{def:rs:event-driven-backtests:fill}{modèle d'exécution} et le portefeuille, tous ne
pouvant programmer que des événements futurs~; et un enregistreur qui produit le résultat et la piste d'audit.
\end{solution}

\begin{solution}{iq:rs:event-driven-backtests:2}
\omterm{def:rs:event-driven-backtests:fill}{Exécutions au contact} (optimistes~: en tête de file), \omterm{def:rs:event-driven-backtests:fill}{exécutions par pénétration} (pessimistes~: en queue de file), exécutions plafonnées au volume (une part du volume de la barre),
exécutions probabilistes étalonnées sur des données réelles~; chacune exige une politique pour les ordres actifs pendant une partie de la barre. Biais~: le contact surestime les exécutions et sous-estime
la sélection défavorable~; la pénétration sous-estime les exécutions~; les plafonds ignorent où le volume a été échangé.
\end{solution}

\begin{solution}{iq:rs:event-driven-backtests:3}
Ordonner totalement les événements d'ordres (d'abord le temps, puis la nature, puis un numéro de séquence), transmettre explicitement des générateurs initialisés, éviter l'itération sur des conteneurs non ordonnés et tout recours
à l'horloge murale, et tester que deux exécutions produisent des résultats identiques.
\end{solution}

\begin{solution}{iq:rs:event-driven-backtests:4}
La tendance des ordres passifs à être exécutés lorsque le prix est sur le point de bouger contre eux, parce que le flux qui les exécute est informé ou persistant.
Mesurer le mark-out de chaque exécution~: le milieu quelques secondes ou barres plus tard, moins le prix d'exécution, dans le sens de l'exécution, moyenné par \omterm{def:rs:event-driven-backtests:fill}{modèle d'exécution}, instant et taille.
\end{solution}

\begin{solution}{iq:rs:event-driven-backtests:5}
Comme un événement à son instant d'effet~: multiplier les positions et les quantités en attente par le ratio, diviser les prix limites et les prix stop par celui-ci, conserver les barres telles qu'elles ont été échangées, et consigner
l'ajustement~; de même pour les regroupements d'actions, avec un contrôle que le capital n'est pas modifié par l'événement.
\end{solution}

\begin{solution}{iq:rs:event-driven-backtests:6}
Lorsque le résultat dépend d'hypothèses que les barres ne peuvent trancher~: le relancer avec des \omterm{def:rs:event-driven-backtests:fill}{exécutions au contact} et par pénétration, et avec des politiques conservatrice et optimiste pour les ordres actifs dans la barre~;
si les réponses encadrent zéro ou diffèrent de plus que l'avantage, il faut recourir au niveau~3 (file et latence) ou disposer de preuves réelles.
\end{solution}
